\scanpage{6}
\mspage{- 6 -}
\begin{flushright}
n\textsuperscript{o}\handwritten{354}
\end{flushright}

\uline{dués de degré} 0. Mais on notera qu'un homomorphisme gradué bijectif\handwritten{\unsure{'}} de\handwritten{\unsure{''}} degré $\ne 0$ n'est pas un isomorphisme de $A$-modules gradués si $M$ et $N$ ne sont pas réduits à 0.

\uline{Exemples}.- 9) Si $M$ est un $A$-module gradué, $M(\lambda_0)$ un $A$-module gradué obtenu par décalage (Exemple 6), l'application identique de $M(\lambda_0)$ sur $M$ est un homomorphisme gradué bijectif de degré $\lambda_0$.

10) Si $a$ est un élément homogène de degré $\delta$ appartenant au centre de $A$, l'endomorphisme $x\to ax$ d'un $A$-module gradué quelconque $M$ est un homomorphisme gradué de degré $\delta$.

2. \uline{Sous-modules gradués.}

PROPOSITION 2.- \uline{Soit} $\Delta$ \uline{un groupe commutatif}. \uline{Soient} $A$ \uline{un anneau gradué de type} $\Delta$, $M$ \uline{un} $A$-\uline{module gradué de type} $\Delta$, $(M_\lambda)$ \uline{sa graduation}, $N$ \uline{un sous}-$A$-\uline{module de} $M$. \uline{Les propriétés suivantes sont équivalentes} :

a) $N$ \uline{est somme des} $N\cap M_\lambda$ $(\lambda\in\Delta)$.

b) \uline{Les composantes homogènes de tout élément de} $N$ \uline{appartiennent à} $N$.

c) $N$ \uline{est engendré par des éléments homogènes}.

Tout élément de $\unsure{N|M}$ s'écrivant d'une seule manière comme somme d'éléments des $M_\lambda$, il est immédiat que a) et b) sont équivalentes, et que a) entraîne c). Montrons que c) implique b). Soit donc $(x_{\handwritten{\unsure{\iota|\tau}}})_{\handwritten{\iota}\in I}$ une famille de générateurs homogènes $\ne 0$ de $N$ et soit $\delta(\handwritten{\iota})$ le degré de $x_{\handwritten{\iota}}$. Tout élément $x\in N$ s'écrit $\sum\limits_{\handwritten{\iota}}a_{\handwritten{\iota}}x_{\handwritten{\iota}}$ avec $a_{\handwritten{\iota}}\in A$ ; si $a_{\handwritten{\iota},\lambda}$ est la composante de degré $\lambda$ de $a_{\handwritten{\iota}}$ $(\lambda\in\Delta)$, on peut donc écrire $x=\sum\limits_{\lambda\in\Delta}\bigl(\sum\limits_{\handwritten{\iota}\in I}a_{\handwritten{\iota},\lambda-\delta(\handwritten{\iota})}x_{\handwritten{\iota}}\bigr)$ ; comme $\sum\limits_{\handwritten{\iota}\in I}a_{\handwritten{\iota},\lambda-\delta(\handwritten{\iota})}x_{\handwritten{\iota}}$ est homogène de degré $\lambda$ et appartient à $N$, on voit que les composantes homogènes de $x$ appartiennent à $N$.

Lorsqu'un sous-module $N$ de $M$ vérifie les propriétés équivalentes énoncées dans la prop.2, il est clair que les $N\cap M_\lambda$ forment une \uline{graduation} compatible avec la structure de module de $N$ ; on dit que $N$, muni de cette graduation,
% archive sheet number: 5

